\documentclass[11pt,a4paper]{article}

\usepackage[margin=2.5cm]{geometry}
\usepackage[utf8]{inputenc}
\usepackage[T1]{fontenc}
\usepackage{amsmath}
\usepackage{hyperref}
\usepackage{cite}

\title{Inteligencia Artificial Explicable y Confiable\\
Fundamentada en Modelos Físicos Analíticos\\[0.5em]
\large Borrador de Introducción -- Propuesta de Investigación\\
Convocatoria Grupos de Investigación Uniovi 2026--2027}

\author{Universidad de Oviedo}
\date{\today}

\begin{document}

\maketitle

\section{Introducción}

La adopción de sistemas de aprendizaje profundo en ámbitos de alta responsabilidad
como la sanidad, las finanzas, la infraestructura crítica y la administración pública
ha situado la transparencia y la confianza entre los requisitos centrales de estos sistemas,
al margen de su rendimiento predictivo \cite{haufe2026explainable, eu_ai_act_2024}.
La naturaleza opaca de estos modelos dificulta que los operadores puedan validar,
auditar y asumir responsabilidad sobre sus decisiones \cite{binkyte2026trustworthy, qi2026towards}.
A esta dificultad técnica se añade una dimensión regulatoria: el Reglamento General
de Protección de Datos (RGPD) y el Reglamento (UE) 2024/1689 de Inteligencia Artificial
establecen obligaciones concretas de explicabilidad y trazabilidad para sistemas
clasificados como de alto riesgo \cite{eu_ai_act_2024}.

La Inteligencia Artificial Explicable (XAI) aborda esta opacidad mediante métodos
que pueden clasificarse según tres criterios: el momento de aplicación
(post-hoc sobre modelos ya entrenados, o ante-hoc con interpretabilidad intrínseca),
el alcance (local para decisiones individuales, o global para el comportamiento general
del modelo) y la dependencia arquitectónica (métodos agnósticos o específicos
de una arquitectura) \cite{lyberatos2026xai, palikhe2025towards}.
Entre los métodos de referencia se encuentran la Propagación de Relevancia por Capas
(LRP) \cite{bach2015pixel}, las Explicaciones Locales Interpretables Agnósticas del
Modelo (LIME) \cite{ribeiro2016trust}, las Explicaciones Aditivas de Shapley
(SHAP) \cite{lundberg2017unified}, la Localización Visual por Gradientes Ponderados
por Clase (Grad-CAM, del inglés \textit{Gradient-weighted Class Activation Mapping})
\cite{selvaraju2017grad} y los Gradientes Integrados \cite{sundararajan2017axiomatic}.
Trabajos posteriores han cuestionado la fiabilidad de estos métodos:
Adebayo et al.\ mostraron que varios de ellos son independientes
de los parámetros del modelo entrenado \cite{adebayo2018sanity}, y Ghorbani et al.\
demostraron que perturbaciones imperceptibles en la entrada pueden modificar sustancialmente
los mapas de saliencias sin alterar la predicción \cite{ghorbani2019interpretation}.

En paralelo, el Grupo de Expertos de Alto Nivel sobre IA de la Comisión Europea
formalizó el concepto de Inteligencia Artificial de Confianza (\textit{Trustworthy AI})
en torno a siete dimensiones: supervisión y control humano, robustez técnica y
seguridad, privacidad y gobernanza de datos, transparencia y explicabilidad,
diversidad y equidad, bienestar social y medioambiental, y rendición de cuentas
\cite{eu_ai_act_2024}.
La explicabilidad es transversal a las demás dimensiones, ya que sin ella no es posible
verificar el cumplimiento de las restantes \cite{binkyte2026trustworthy, binkyte2026causality}.
La Ley de IA de la UE (Reglamento 2024/1689) concreta estas obligaciones:
el Artículo 86 reconoce el derecho individual a recibir una explicación de las
decisiones de sistemas de alto riesgo; los Anexos~I y III delimitan los ámbitos
de aplicación, entre ellos dispositivos médicos, infraestructura crítica,
evaluación crediticia, selección de empleo y biometría,
con entrada en vigor el 2~de agosto de 2026 \cite{eu_ai_act_2024}.

\subsection{Limitaciones actuales de XAI}

La investigación reciente ha identificado una debilidad de base en los métodos de XAI:
el diseño \textit{algorithm-first} \cite{haufe2026explainable}.
Haufe et al.\ argumentan que los métodos dominantes, incluidos SHAP y Grad-CAM,
se desarrollaron como algoritmos sin una definición previa del problema que deben resolver
ni criterios objetivos de corrección \cite{haufe2026explainable}.
Como resultado, estos métodos pueden atribuir importancia a variables que están
correlacionadas con el objetivo pero que no son causalmente relevantes,
generando atribuciones que los operadores interpretan como válidas
sin que exista garantía de que lo sean
\cite{haufe2026explainable, adebayo2018sanity}.

Un segundo problema es la ausencia de verdad de terreno (\textit{ground-truth}) objetiva
para evaluar la corrección de las explicaciones.
Haufe et al.\ señalan que no existe ground-truth para el problema de XAI en conjuntos
de datos reales, lo que obliga a recurrir a sustitutos como el juicio humano o métricas
de fidelidad que no garantizan la corrección de la explicación
\cite{haufe2026explainable}.
Palikhe et al.\ documentan la misma limitación en el ámbito de los modelos de lenguaje,
donde las verdades de terreno explicativas son con frecuencia inaccesibles \cite{palikhe2025towards}.
Gipi\v{s}kis y Kurasova proponen las Tarjetas de Evaluación de XAI como mecanismo
de estandarización, que exigen declarar el contexto de validez y el riesgo de
manipulación de cada métrica \cite{gipiskis2026evaluation}.
Esta situación guarda paralelismo con la saturación que afecta a los
\textit{benchmarks} generales de IA: en el primer semestre de 2026, evaluaciones como
MMLU-Pro o GPQA Diamond presentan diferencias entre modelos dentro del margen de ruido
estadístico, lo que reduce su utilidad discriminativa
\cite{kili_benchmarks_2026, stanford_hai_index_2026}.

Los modelos de aprendizaje profundo tienden a capturar correlaciones presentes en los
datos de entrenamiento que no corresponden a mecanismos causales del fenómeno estudiado
\cite{binkyte2026causality, ghorbani2019interpretation}.
En el análisis de electroencefalografía (EEG), por ejemplo, los clasificadores de
crisis epilépticas pueden aprender a correlacionarse con artefactos musculares o de
movimiento ocular en lugar de con los patrones de actividad ictal \cite{lyberatos2026xai}.
Los métodos post-hoc actuales no pueden distinguir estas correlaciones espurias de
las causalmente válidas sin una referencia externa al dominio:
la importancia que SHAP o Grad-CAM asignan a una variable no informa sobre si esa
importancia es causalmente justificada o un artefacto del entrenamiento
\cite{haufe2026explainable, lyberatos2026xai}.

Desde la perspectiva regulatoria, ninguno de los enfoques actuales de XAI ofrece
certificabilidad formal de sus explicaciones.
Los Modelos Causales Estructurales (SCM) constituyen una línea de trabajo relevante:
Binkyte et al.\ los proponen como mecanismo para resolver los conflictos entre las
dimensiones de IA de Confianza \cite{binkyte2026trustworthy, binkyte2026causality}.
Sin embargo, los SCM son estructuras aprendidas de datos y pueden heredar sus sesgos,
por lo que no garantizan validez fuera de la distribución de entrenamiento
\cite{binkyte2026causality}.
El Artículo 86 y los Anexos~I y III de la Ley de IA de la UE exigen trazabilidad
documental y evaluación de conformidad que los enfoques estadísticos no pueden
satisfacer de forma verificable e independiente de los datos \cite{eu_ai_act_2024}.
Qi et al.\ documentan que solo el 30\% de las organizaciones han establecido controles
de gobernanza adaptados a sistemas autónomos, lo que refleja en parte la ausencia
de mecanismos formales de auditoría \cite{qi2026towards}.

\subsection{Uso de modelos físicos analíticos en XAI}

Este trabajo parte de la observación de que en numerosos dominios de aplicación
existen modelos físicos analíticos bien establecidos: ecuaciones de primer principio,
modelos de estado, modelos farmacocinéticos y farmacodinámicos (PK/PD) y modelos dipolares
neurofisiológicos, entre otros.
Estos modelos codifican relaciones causales del dominio con independencia de los datos
de entrenamiento, y su validez no está restringida a una distribución observada
\cite{lyberatos2026xai, binkyte2026causality}.
Lyberatos et al.\ identifican la incorporación de \textit{priors} de dominio,
como la adyacencia espacial de electrodos o restricciones espectrales de origen
neurofisiológico, como una dirección necesaria para dotar de fundamento las
explicaciones en EEG, aunque sin un marco formal que lo generalice \cite{lyberatos2026xai}.
La propuesta de este trabajo es formalizar ese rol: el modelo físico analítico
define qué atribuciones son admisibles y permite detectar aquellas que contradigan
las relaciones causales del dominio.

Una segunda aportación del uso de modelos físicos es la generación de
verdad de terreno mediante simulación.
El modelo analítico puede producir escenarios con resultado conocido por construcción
a partir de las ecuaciones del dominio, sin depender de datos etiquetados,
lo que permite medir la tasa de error del explicador de forma objetiva y reproducible
\cite{haufe2026explainable, qi2026towards}.
Los sistemas agénticos emplean un principio similar mediante \textit{world models}
para verificar las consecuencias de acciones antes de ejecutarlas \cite{qi2026towards},
pero este mecanismo no se ha trasladado a la evaluación de explicaciones en sistemas
de aprendizaje supervisado.
El protocolo de certificación resultante es compatible con los requisitos de
trazabilidad documental del Art.~86 y los Anexos~I y III de la Ley de IA
\cite{eu_ai_act_2024}.

Una tercera función es la detección de correlaciones espurias.
Si el explicador asigna importancia elevada a una variable que el modelo físico indica
como irrelevante para el fenómeno, esa discrepancia señala una posible correlación
espuria aprendida durante el entrenamiento.
Esta capacidad complementa el paradigma de invarianza causal de Binkyte et al.\
\cite{binkyte2026trustworthy}: mientras los SCM identifican correlaciones espurias
por su inestabilidad entre distribuciones, el modelo físico las identifica por
contradicción directa con las leyes del dominio, con independencia de los datos
\cite{binkyte2026causality, lyberatos2026xai}.
Que la detección no dependa de los datos de entrenamiento es relevante para la
certificación formal, ya que evita la circularidad de validar el modelo con los
mismos datos con que fue entrenado \cite{eu_ai_act_2024}.

Este trabajo propone el marco de \textbf{Inteligencia Artificial Explicable y Confiable
Fundamentada en Modelos Físicos Analíticos} (PG-XAI, del inglés
\textit{Physics-Grounded XAI}), que integra modelos físicos analíticos como
referencia externa de fidelidad para explicadores de aprendizaje automático
en entornos regulados.
El objetivo es desarrollar y evaluar los tres mecanismos descritos en
\textbf{[CASOS DE APLICACIÓN]}.
La Sección~2 revisa el estado del arte en XAI y modelos físicos en los
dominios considerados; la Sección~3 formaliza el marco PG-XAI;
la Sección~4 describe el diseño experimental; la Sección~5 presenta
el plan de trabajo para el periodo 2026--2027.

\bibliographystyle{unsrt}
\bibliography{ref}

\end{document}
